\clearpage
\phantomsection
\section*{Chương 7. Tự tương quan --- MAC mặc định / FFT tùy chọn + BFP}
\addcontentsline{toc}{section}{Chương 7. Tự tương quan — MAC mặc định / FFT tùy chọn + BFP}
\pdfbookmark[1]{Chương 7. Tự tương quan — MAC mặc định / FFT tùy chọn + BFP}{chuong7}

\begin{learningobj}
Sau chương này bạn viết được định nghĩa tự tương quan không vòng từ nguyên lý đầu tiên; tính tay từng giá trị $\mathcal{R}[\tau]$ trên chuỗi ngắn và trên khung cosin chu kỳ~8; thiết kế datapath MAC tuần tự bằng một khối DSP và bộ tích lũy rộng kèm bảng timing chu kỳ; giải thích FFT IP radix-4 / BFP từng tầng cánh bướm và barrel shift $2^{2e}$ trước IFFT; viết twin Python bit-exact; và tính được ngân sách chu kỳ / RTF so với bước nhảy $10\,\text{ms}$.
\end{learningobj}

\subsection*{7.1. Hiện tượng thật: đo mức ``giống chính mình''}

Khi dây thanh quản rung gần tuần hoàn với chu kỳ cơ bản $T_0$ giây, sóng áp suất tại thời điểm $t$ gần trùng với sóng tại $t+T_0$. Trên lưới mẫu $f_s = 16\,\text{kHz}$, chu kỳ đó là $\tau_0 = f_s\cdot T_0$ mẫu. Ví dụ giọng nam khoảng $100\,\text{Hz}$ cho $\tau_0 = 160$; giọng nữ khoảng $200\,\text{Hz}$ cho $\tau_0 = 80$. Phép \textbf{tự tương quan (autocorrelation)} đo mức giống nhau đó bằng cách trượt một bản sao đi $\tau$ mẫu (\textbf{độ trễ --- lag}), nhân từng cặp mẫu chồng nhau và cộng dồn. Tại $\tau = \tau_0$, các cặp cùng pha cho tích dương lớn, tạo đỉnh nổi trên đồ thị $\mathcal{R}_x[\tau]$. Đo được $\tau_0$ thì suy ra ngay $F_0 = f_s/\tau_0$.

Đây là hiện tượng vật lý thật của sóng tuần hoàn, không chỉ là quy ước toán học. Định lý \textbf{Wiener--Khinchin} (Chương~0) chứng minh rằng $\mathcal{R}_x$ cũng bằng biến đổi Fourier ngược của mật độ phổ công suất $|X[k]|^2$, nên có hai đường silicon tương đương về toán: nhân--cộng trực tiếp (MAC) hoặc FFT $\to |X|^2 \to$ IFFT. Với dải trễ hẹp $\tau = 26 \ldots 214$ (189 giá trị, khớp \texttt{lab\_1\_3\_5.py} và bản đồ Chương~0), đường mặc định trên Cyclone~V là \textbf{MAC} (\emph{Multiply-ACcumulate}): chỉ tính đúng những $\tau$ cần thiết, không tính toàn bộ phổ rồi mới cắt. Mỗi hop $10\,\text{ms}$ chỉ có $500\,000$ chu kỳ ở $50\,\text{MHz}$; $98\,280$ phép MAC chiếm khoảng một phần năm ngân sách đó và vẫn để lại biên cho cửa sổ, chuẩn hóa và chia $F_0$.

Tại sao không lấy đỉnh phổ trực tiếp? Vì phổ biên độ của nguyên âm có đỉnh formant (cộng hưởng ống thanh) có thể cao hơn đỉnh $F_0$; tự tương quan thời gian ``nhìn'' chu kỳ lặp của toàn sóng, ít bị formant đánh lừa hơn --- đó là lý do Praat chọn ACF (sau chuẩn hóa) làm lõi đo cao độ.

\begin{physicalinsight}
Tự tương quan không ``tạo ra'' cao độ. Nó chỉ đo độ trùng pha giữa tín hiệu và bản sao trễ của chính nó. Nếu khung đã bị rò phổ hoặc còn DC (Chương~3,~6), đỉnh giả sẽ cạnh tranh với đỉnh thật --- MAC tính đúng công thức vẫn báo sai $F_0$.
\end{physicalinsight}

\begin{figure}[H]
\centering
\begin{tikzpicture}[font=\small, >=Stealth]
  \draw[->] (0,0)--(9.0,0) node[right]{$\tau$ (mẫu)};
  \draw[->] (0,0)--(0,3.0) node[above]{$\mathcal{R}_x[\tau]$};
  \draw[thick, teal] (0.3,2.7)--(1.2,0.9)--(2.4,2.2)--(3.6,0.7)--(4.8,1.6)--(6.0,0.5)--(7.2,1.1)--(8.4,0.4);
  \fill[red!70] (2.4,2.2) circle (2.2pt) node[above right,font=\scriptsize]{$\tau_0=T_0$};
  \fill[orange!80] (4.8,1.6) circle (2.2pt) node[above,font=\scriptsize]{$2T_0$};
  \node[font=\scriptsize, align=left] at (7.0,2.4){đỉnh phụ bội\\có thể gây sai quãng tám};
\end{tikzpicture}
\caption{Hàm tự tương quan của khung hữu thanh: đỉnh toàn cục tại $\tau=0$ (năng lượng), đỉnh cục bộ đầu tiên tại $\tau_0=T_0$ là ứng viên cao độ; đỉnh $2T_0$ là ứng viên sai quãng tám.}
\label{fig:ch07_acf_peaks}
\end{figure}

\subsection*{7.2. Toán từng bước --- định nghĩa gốc trước mọi phím tắt}

Hàm tự tương quan \textbf{không vòng (linear / non-circular)} của khung $x[n]$ độ dài $L$ được định nghĩa từ tích vô hướng của hai đoạn chồng:
\[
\mathcal{R}_x[\tau] = \sum_{n=0}^{L-1-\tau} x[n]\cdot x[n+\tau], \qquad \tau \ge 0.
\]
Từng bước: (1)~chọn $\tau$; (2)~lấy cặp $(x[0],x[\tau])$, $(x[1],x[\tau+1])$, \ldots, $(x[L-1-\tau],x[L-1])$; (3)~nhân từng cặp; (4)~cộng tất cả tích. Không được viết tắt thành ``chỉ một hạng $x[0]x[\tau]$''.

Mỗi ký hiệu:
\begin{itemize}[leftmargin=1.4em]
\item $x[n]$: mẫu đã trừ mean và đã nhân cửa sổ Hanning (Chương~3,~6), định dạng Q1.15 trên silicon.
\item $\tau$: độ trễ tính bằng số mẫu. Dải quét phần cứng: $\tau \in [26,214]$ (189 giá trị) để đỉnh tìm trong $[27,213]$ luôn có láng giềng cho nội suy parabolic (Chương~9). Sàn $75\,\text{Hz}$ $\leftrightarrow$ $\tau\approx 213$; trần $600\,\text{Hz}$ $\leftrightarrow$ $\tau\approx 27$.
\item Số hạng trong tổng đúng bằng $L-\tau$: trễ càng lớn, càng ít cặp chồng --- đây là thiên vị hình học sẽ được bù ở Chương~8.
\item Tại $\tau=0$: $\mathcal{R}_x[0] = \sum_n x[n]^2$ là năng lượng khung (cực đại toàn cục).
\end{itemize}

Định lý Wiener--Khinchin cho phép đường tương đương:
\[
\mathcal{R}_x[\tau] = \mathcal{F}^{-1}\{|X[k]|^2\}.
\]
FFT mặc định tính tự tương quan \textbf{vòng (circular)}: mẫu cuối ``quấn'' sang mẫu đầu. Muốn khớp định nghĩa không vòng ở trên phải zero-pad ít nhất thêm $L$ số 0 trước FFT (thường pad lên $N_{\mathrm{FFT}}=1024$ hoặc $2048$). Đó là một trong những lý do đường mặc định pitch chọn MAC.

Trên silicon, mỗi tích $x[n]\cdot x[n+\tau]$ là Q1.15$\times$Q1.15 $=$ Q2.30 (32-bit). Cộng dồn đến $640$ hạng tử đòi hỏi bộ tích lũy rộng: thiết kế dùng \textbf{41-bit có dấu} (khớp stub \texttt{rtl/mac\_autocorr.v} và comment trong \texttt{lab\_1\_3\_5.py}). Ước lượng thô: $|x|\le 2^{15}$, $|x\cdot x|\le 2^{30}$, $640\cdot 2^{30} < 2^{40}$ --- 41 bit có dấu vừa đủ biên an toàn một bit.

\subsection*{7.3. Ví dụ tính tay --- mọi số đã chạy Python}

\paragraph{A. Chuỗi cực ngắn $x = [1,2,1,2]$.}
Từng cặp:
\begin{align*}
\mathcal{R}[0] &= 1\cdot 1 + 2\cdot 2 + 1\cdot 1 + 2\cdot 2 = 1+4+1+4 = 10,\\
\mathcal{R}[1] &= 1\cdot 2 + 2\cdot 1 + 1\cdot 2 = 2+2+2 = 6,\\
\mathcal{R}[2] &= 1\cdot 1 + 2\cdot 2 = 1+4 = 5,\\
\mathcal{R}[3] &= 1\cdot 2 = 2.
\end{align*}
Tỉ số thô $\mathcal{R}[2]/\mathcal{R}[0] = 0{,}5$. Sau khi nhân $C[\tau]$ (Chương~8) và chia theo số hạng, đỉnh chu kỳ được kéo về gần 1.

\paragraph{B. Khung cosin chu kỳ 8 mẫu, $N=16$ (Q1.15).}
$x[n] = \mathrm{sat}_{16}(\mathrm{round}(\cos(2\pi n/8)\cdot 32768))$, $n=0..15$. Từng bước cho $n=0$: $\cos 0 = 1$, $\mathrm{round}(32768)=32768$, sat $\to 32767$. Cho $n=4$: $\cos\pi = -1$, $\mathrm{round}(-32768)=-32768$ (vừa khít min int16).
\begin{center}\footnotesize
\begin{tabular}{@{}lllllllll@{}}\toprule
$n$ & 0 & 1 & 2 & 3 & 4 & 5 & 6 & 7 \\
$x$ & 32767 & 23170 & 0 & $-23170$ & $-32768$ & $-23170$ & 0 & 23170 \\
\midrule
$n$ & 8 & 9 & 10 & 11 & 12 & 13 & 14 & 15 \\
$x$ & 32767 & 23170 & 0 & $-23170$ & $-32768$ & $-23170$ & 0 & 23170 \\
\bottomrule
\end{tabular}
\end{center}
Tự tương quan không vòng (script):
\begin{align*}
\mathcal{R}[0] &= 8\,589\,627\,426,\\
\mathcal{R}[8] &= 4\,294\,813\,713 \approx \tfrac12\mathcal{R}[0],\\
\mathcal{R}[4] &= -6\,442\,220\,568 \quad\text{(chống pha, không phải đỉnh dương)}.
\end{align*}
Tỉ số thô $\mathcal{R}[8]/\mathcal{R}[0] = 0{,}5$ vì chỉ còn 8 cặp; sau chuẩn hóa unbiased $(\mathcal{R}[8]/8)/(\mathcal{R}[0]/16) = 1$. Đây là lý do Chương~8 tồn tại.

\paragraph{C. Ngân sách chu kỳ trên đường MAC thật.}
Với $\tau = 26\ldots 214$, tổng số phép nhân--cộng đúng bằng
\[
N_{\mathrm{MAC}} = \sum_{\tau=26}^{214}(640-\tau).
\]
Tính tổng cấp số: số hạng $189$; số hạng đầu $640-26=614$; số hạng cuối $640-214=426$; tổng $= 189\cdot(614+426)/2 = 98\,280$.
Không phải $189\times 640 = 120\,960$ --- công thức naïve đếm thừa. Tại $f_{\mathrm{clk}} = 50\,\text{MHz}$:
\[
T_{\mathrm{compute}} = \frac{98\,280}{50\times 10^6} = 1{,}9656\,\text{ms}, \qquad
\mathrm{RTF} = \frac{1{,}9656}{10} = 0{,}19656 \ll 1.
\]

\subsection*{7.4. Ý nghĩa phần cứng --- MAC mặc định}

Hai con trỏ đọc khung trong BRAM True Dual-Port (Chương~5): một trỏ tới $x[n]$, một tới $x[n+\tau]$. Mỗi chu kỳ một khối DSP nhân $18\times 18$ (hoặc $18\times 19$), cộng vào bộ tích lũy 41-bit có dấu. Quét $\tau$ tuần tự. Sau $(L-\tau)$ bước, thanh ghi accum chính là $\mathcal{R}_x[\tau]$.

\begin{figure}[H]
\centering
\begin{tikzpicture}[font=\footnotesize, >=Stealth]
  \node[draw, rounded corners] (a) {RAM A $x[n]$};
  \node[draw, rounded corners, below=0.55cm of a] (b) {RAM B $x[n{+}\tau]$};
  \node[draw, rounded corners, fill=green!12, right=1.6cm of a] (m) {DSP $\times$};
  \node[draw, rounded corners, fill=yellow!15, right=1.3cm of m] (acc) {ACC 41-bit};
  \draw[->] (a) -| (m); \draw[->] (b) -| (m); \draw[->] (m) -- (acc);
  \node[draw, rounded corners, below=1.1cm of m, align=center] (rom) {$C[\tau]$ ROM\\(Ch.~8)};
  \node[draw, rounded corners, right=1.1cm of rom] (s) {$S[\tau]$};
  \draw[->] (acc) |- (s); \draw[->] (rom) -- (s);
\end{tikzpicture}
\caption{Datapath MAC: hai toán hạng từ TDP RAM, một nhân DSP, một tích lũy --- khớp vòng \texttt{for tau} trong \texttt{lab\_1\_3\_5.py}.}
\label{fig:ch07_mac_datapath}
\end{figure}

\paragraph{Bảng timing chu kỳ --- một $\tau$ cố định (ví dụ $\tau=2$, $L=4$).}
\begin{center}\small
\begin{tabular}{@{}cllll@{}}\toprule
ck & Trạng thái & $n$ & $xa{=}x[n]$, $xb{=}x[n{+}\tau]$ & \texttt{acc} sau cạnh \\
\midrule
0 & IDLE$\to$LOAD & --- & --- & clear $0$ \\
1 & RUN & 0 & $x[0],x[2]$ & $x[0]x[2]$ \\
2 & RUN & 1 & $x[1],x[3]$ & $x[0]x[2]+x[1]x[3]$ \\
3 & DONE & --- & --- & $=\mathcal{R}[2]$; \texttt{done\_tau}=1 \\
\bottomrule
\end{tabular}
\end{center}
Với $\tau=26$, $L=640$: pha RUN kéo dài đúng $614$ chu kỳ. FSM ngoài tăng $\tau$ và lặp. Không được clear \texttt{acc} giữa các $n$ của cùng một $\tau$.

\paragraph{RTL lõi (ý tưởng, khớp stub \texttt{mac\_autocorr.v}).}
\begin{verbatim}
module mac_autocorr (
  input  wire               clk, rst_n, start,
  input  wire signed [15:0] xa, xb,   // x[n], x[n+tau] Q1.15
  input  wire [9:0]         n_pairs,  // = L - tau
  output reg  signed [40:0] acc,      // 41-bit signed
  output reg                done_tau
);
  // FSM: IDLE -> RUN (n_pairs cycles: acc <= acc + xa*xb) -> DONE
  // xa, xb phải ổn định đúng cạnh khi acc cập nhật (đồng bộ với
  // latency đọc TDP RAM --- thường 1 ck registered).
endmodule
\end{verbatim}
Không dùng \texttt{numpy.correlate} float làm chuẩn silicon: twin phải nhân--cộng int cùng độ rộng accum.

\subsection*{7.5. Đường tùy chọn: FFT IP + BFP + barrel shift trước IFFT}

Chỉ bật sau khi MAC đã xanh bit-exact. Dùng \textbf{FFT Intel FPGA IP} (User Guide~683374):

\begin{itemize}[leftmargin=1.4em]
\item Kiến trúc \textbf{radix-4} (hoặc radix-2): mỗi tầng cánh bướm (\textbf{butterfly}) kết hợp 4 (hoặc 2) điểm với twiddle $W_N^k = e^{-j2\pi k/N}$ lấy từ ROM nội bộ. Với $N=1024$ radix-4 cần $\log_4 1024 = 5$ tầng; radix-2 cần 10 tầng.
\item Streaming hoặc buffered; tín hiệu SOP/EOP đánh dấu một block $N$ điểm trên Avalon-ST.
\item \textbf{Block Floating Point (BFP --- số chấm động khối)}: mọi bin chia sẻ một số mũ nguyên $e$; mantissa dịch sao cho $\max|m|$ vừa khít bit. Sau mỗi tầng butterfly, nếu có nguy cơ tràn thì toàn vector dịch phải 1 bit và $e$ tăng 1. IP xuất \texttt{source\_exp}.
\item BFP \textbf{không} có cổng nhập exponent vào IFFT --- kỹ sư phải scale bên ngoài trước khi đưa mantissa vào IFFT.
\end{itemize}

\paragraph{BFP từng tầng --- ví dụ mantissa 8-bit, 2 tầng.}
Bắt đầu stage~0: $m = [40,-30,10,-5]$, $e = 0$. Ngưỡng tràn minh họa $= 32$ (giả sử còn 1 bit đầu). Vì $\max|m| = 40 > 32$:
\begin{center}\small
\begin{tabular}{@{}clll@{}}\toprule
Sự kiện & $e$ & Mantissa & Ghi chú \\
\midrule
Vào stage 0 & 0 & $[40,-30,10,-5]$ & $\max=40>32$ \\
Sau scale stage 0 & 1 & $[20,-15,5,-2]$ & dịch phải 1 \\
Sau stage 1 (không tràn) & 1 & $[20,-15,5,-2]$ & giữ $e$ \\
Power bin 0 thô $|m|^2$ & --- & $400$ & quên $2^{2e}$ \\
Power đúng $|m|^2\cdot 2^{2e}$ & --- & $1600$ & $=40^2$ \\
\bottomrule
\end{tabular}
\end{center}

\begin{firstprin}{Sau $|X|^2$, số mũ nhân đôi thành $2e$}
Mantissa $y = 3$, exponent $e = 2$ có giá trị thật $3\cdot 2^{-2}$. Bình phương:
\[
(3\cdot 2^{-2})^2 = 9\cdot 2^{-4}
\]
--- số mũ thành $2e$. Đẩy mantissa $|y|^2$ thẳng vào IFFT (IP kỳ vọng fixed-point với exponent $=0$ ở lối vào) mà quên barrel-shift theo $2e$ sẽ làm cả dải động $\mathcal{R}[\tau]$ sai tỷ lệ, phá ngưỡng voicing.
\end{firstprin}

\paragraph{Barrel shift trước IFFT.}
RTL bắt buộc: nếu \texttt{source\_exp}$=e$ sau FFT, sau khi tính $|X[k]|^2$ trên mantissa, dịch trái (hoặc nhân) kết quả theo $2e$ trước khi đưa vào IFFT --- hoặc dịch mantissa IFFT input tương đương. Twin phải dùng model \texttt{\_model} của variation FFT IP, không đoán $e$.

\begin{figure}[H]
\centering
\begin{tikzpicture}[font=\scriptsize, >=Stealth]
  \node[draw] (x) {$x$};
  \node[draw, right=0.5cm of x] (fft) {FFT};
  \node[draw, fill=orange!12, right=0.5cm of fft] (sh1) {theo $e$};
  \node[draw, right=0.5cm of sh1] (mag) {$|X|^2$};
  \node[draw, fill=orange!12, right=0.5cm of mag] (sh2) {barrel $2e$};
  \node[draw, right=0.5cm of sh2] (ifft) {IFFT};
  \node[draw, right=0.5cm of ifft] (r) {$\mathcal{R}$};
  \draw[->] (x)--(fft)--(sh1)--(mag)--(sh2)--(ifft)--(r);
  \node[draw, dashed, below=0.75cm of fft, fill=gray!8] (exp) {\texttt{source\_exp}$=e$};
  \draw[->, dashed] (fft) -- (exp);
  \draw[->, dashed] (exp.east) -| (sh1.south);
\end{tikzpicture}
\caption{Đường FFT/Wiener--Khinchin: hai khối cam (scale theo $e$ và barrel $2^{2e}$) là RTL bắt buộc. Quên chúng gần như chắc fail bit-exact --- đó là lý do mặc định vẫn là MAC.}
\label{fig:ch07_fft_bfp}
\end{figure}

\paragraph{Vì sao MAC thắng khi $\tau_{\max}\ll L$.}
MAC: $98\,280$ nhân thực/khung. FFT 1024 radix-4: 5 tầng, mỗi tầng $256$ butterfly phức (mỗi butterfly vài nhân thực) --- cộng IFFT tương tự cộng $|X|^2$, rồi vẫn phải cắt $\tau$. Thêm rủi ro bit-reverse và BFP. Chỉ mở FFT khi lab yêu cầu so khớp đường Praat FFT hoặc khi RTF MAC không đủ.


\paragraph{Bit-reverse và thứ tự bin.}
FFT IP có thể xuất dữ liệu theo thứ tự bit-reverse hoặc natural order tùy cấu hình. Twin và RTL sau $|X|^2$ phải cùng thứ tự trước IFFT. Sai thứ tự bin làm $\mathcal{R}[\tau]$ ``xáo'' --- đỉnh $\tau_0$ có thể biến mất hoặc nhảy chỗ. Lab nghiệm thu nhánh FFT bắt buộc so dump toàn bộ $|X[k]|^2$ với model IP trước khi nối IFFT.

\paragraph{Zero-pad cụ thể.}
Khung $L=640$ pad thêm ít nhất 640 số 0 lên $N_{\mathrm{FFT}}=2048$ (hoặc 1024 nếu chấp nhận circular leak nhỏ và đã đo sai số). Pad ngắn hơn $2L-1$ làm tự tương quan vòng ``đụng'' đuôi vào đầu --- $\mathcal{R}[\tau]$ tại $\tau$ lớn bị cộng thêm hạng giả. MAC không bao giờ gặp lỗi này vì định nghĩa gốc chỉ cộng $L-\tau$ hạng thật.

Trên DE10-Standard, chọn MAC mặc định còn vì một lý do vận hành: khi debug cent-error ở Chương~13, kỹ sư chỉ cần so từng $\mathcal{R}[\tau]$ int với twin; không phải lần theo \texttt{source\_exp} từng block. FFT giữ trong tài liệu như đường tùy chọn để hiểu Wiener--Khinchin và để lab nâng cao so khớp Praat khi bật biến thể phổ.

\subsection*{7.6. Twin Python bit-exact}

\begin{verbatim}
import numpy as np

def acf_linear_int(x, tau_min, tau_max):
    x = np.asarray(x, dtype=np.int64)
    L = len(x)
    out = {}
    for tau in range(tau_min, tau_max + 1):
        # Khớp MAC: cộng dồn int, không float
        out[tau] = int(np.dot(x[: L - tau], x[tau:]))
    return out

# Khớp lab: tau 26..214, accum int64 (>= 41-bit signed trên RTL)
R = acf_linear_int(xw_q15, 26, 214)
# So dump R[tau] với ModelSim: sai lệch 0 bit trên đường MAC
\end{verbatim}
Giải thích: \texttt{np.dot} trên \texttt{int64} tương đương vòng RUN của FSM; không lấy \texttt{numpy.correlate} float. Trên nhánh FFT, twin phải gồm cả barrel $2^{2e}$ giống RTL và so với \texttt{\_model} của IP.


Khi đo RTF thực tế trên board, đừng chỉ đếm chu kỳ MAC lý thuyết: thêm latency đọc TDP RAM (1~ck/mẫu), handshake Avalon-ST có thể chen \texttt{waitrequest}, và pha nhân $C[\tau]$ (Chương~8). Ngân sách an toàn lab: dành $3\,\text{ms}$ cho toàn bộ Pha~2 (cửa sổ + MAC + chuẩn hóa + peak), tương đương RTF $0{,}3$. Nếu sau place-and-route Fmax tụt dưới $40\,\text{MHz}$, phải pipeline thêm nhân DSP hoặc giảm $\tau_{\max}$ --- không được ``hy vọng'' timing tự hết.

\subsection*{7.7. Chi phí silicon}

\begin{costbox}
\textbf{Ngân sách MAC autocorr (default).}
\begin{itemize}[leftmargin=1.2em]
\item \textbf{DSP}: 1 khối MAC/chu kỳ; $98\,280$ phép/khung $\approx 1{,}97\,\text{ms}$ @ $50\,\text{MHz}$.
\item \textbf{ALM}: $\sim 200$ cho FSM con trỏ $\tau$, $n$, handshake Avalon-ST.
\item \textbf{M10K}: đệm khung dùng chung Chương~5; MAC không thêm BRAM.
\item \textbf{RTF} $= 0{,}197 \ll 1$ (còn $>80\%$ thời gian rảnh mỗi hop $10\,\text{ms}$ $= 500\,000$ chu kỳ). Vì sao vừa khít: ngay cả khi thêm cửa sổ + $C[\tau]$ + parabolic + chia, tổng vẫn dưới khoảng $3\,\text{ms}$, RTF $<0{,}3$.
\item \textbf{FFT optional}: hàng nghìn ALM + nhiều M10K twiddle/buffer + DSP theo variation; barrel shifter vài chục ALM. Chỉ bật khi twin đã khóa và lab yêu cầu.
\item Vì sao MAC mặc định: 0 rủi ro BFP/bit-reverse, 0-bit với golden \texttt{lab\_1\_3\_5.py}.
\end{itemize}
\end{costbox}


Tóm lại: đường MAC là hợp đồng bit-exact mặc định của gói này; đường FFT/BFP là bài học kiến trúc và bài lab nâng cao. Mọi số trong mục 7.3 đã chạy qua \texttt{verify/verify\_chuong\_06\_09.py} --- nếu sửa công thức pad hoặc độ rộng accum, sửa script trước rồi mới sửa RTL.

\subsection*{7.8. Bài tập và lời giải giải thích}

\paragraph{Bài tập 7.1.} Với $x = [1,2,1,2]$, tính $\mathcal{R}[1]$ và $\mathcal{R}[2]$ từ định nghĩa gốc. Có được phép viết $\mathcal{R}[2] = x[0]x[2]$ một hạng duy nhất không?

\paragraph{Bài tập 7.2.} Trên khung cosin mục 7.3.B, vì sao $\mathcal{R}[8]/\mathcal{R}[0] = 0{,}5$ dù tín hiệu lặp hoàn hảo? Làm thế nào để tỉ số về gần 1?

\paragraph{Bài tập 7.3.} Mantissa sau FFT $m = 20$, $e = 1$. Tính power $|m|^2\cdot 2^{2e}$. Nếu quên nhân $2^{2e}$, sai bao nhiêu lần?

\paragraph{Bài tập 7.4.} Tính $N_{\mathrm{MAC}}$ đúng bằng $\sum_{\tau=26}^{214}(640-\tau)$ và RTF @ $50\,\text{MHz}$. So với ước lượng naïve $189\times 640$.

\paragraph{Bài tập 7.5.} Vẽ bảng timing 4 chu kỳ cho $\tau=1$, $L=4$, $x=[1,2,1,2]$: mỗi cạnh \texttt{acc} mang giá trị gì? \texttt{done\_tau} lên khi nào?

\noindent\textbf{Lời giải giải thích.}
\begin{enumerate}[leftmargin=1.4em]
\item $\mathcal{R}[1] = 6$, $\mathcal{R}[2] = 5$. Không được chỉ lấy một hạng: định nghĩa cộng mọi cặp chồng, với $\tau=2$ còn đúng hai hạng $1\cdot 1 + 2\cdot 2$.
\item Vì chỉ còn $N-\tau = 8$ cặp trong tổng không vòng, trong khi $\mathcal{R}[0]$ có 16 hạng. Chuẩn hóa unbiased chia cho số hạng, hoặc nhân $C[\tau] = \mathcal{R}_w[0]/\mathcal{R}_w[\tau]$ (Chương~8), kéo tỉ số về gần 1.
\item Power $= 400\cdot 4 = 1600$. Quên $2^{2e}$ còn 400 --- sai đúng $4 = 2^{2e}$ lần (khoảng $+6\,\text{dB}$ nếu đo năng lượng sai chiều ngược), đủ phá ngưỡng voicing.
\item $N_{\mathrm{MAC}} = 98\,280$, $T = 1{,}9656\,\text{ms}$, $\mathrm{RTF} \approx 0{,}197$. Naïve $120\,960$ lớn hơn khoảng $23\%$ vì đếm thừa các hạng không tồn tại khi $\tau>0$.
\item ck1: \texttt{acc}=$1\cdot 2=2$; ck2: $2+2\cdot 1=4$; ck3: $4+1\cdot 2=6=\mathcal{R}[1]$; \texttt{done\_tau} sau đúng $L-\tau=3$ bước RUN.
\end{enumerate}
